\chapter{\MakeUppercase{Implementation of the System}}

This chapter details the various modules and submodules that constitute the implementation of the proposed intelligent surveillance system. Each section describes the methodologies, algorithms, and technologies employed in realizing the functionalities of the system, from initial data acquisition to final output and user interaction.

\section{\MakeUppercase{Frame Extraction}}
This initial module extracts individual frames from a video input, which may be a pre-recorded file or a live stream. The process is implemented using OpenCV (\texttt{cv2}), a Python library for efficient video processing.

\begin{algorithm}[H]
\caption{Frame Extraction from Video Source}
\label{alg:frame_extraction}
\textbf{Input:} Video source (file path or camera index)\\
\textbf{Output:} Sequence of extracted video frames
\begin{algorithmic}[1]
\Function{ExtractFrames}{source}
    \State $video \gets \text{cv2.VideoCapture}(source)$
    \State $frames \gets [\hspace{0.2em}]$
    \If{$video.isOpened(\hspace{0.1em})$}
        \While{True}
            \State $ret, frame \gets video.read(\hspace{0.1em})$
            \If{$\neg ret$}
                \State \textbf{break}
            \EndIf
            \State Append $frame$ to $frames$
        \EndWhile
        \State $video.release(\hspace{0.1em})$
    \Else
        \State \textbf{Error:} Could not open video source
    \EndIf
    \State \Return $frames$
\EndFunction
\end{algorithmic}
\end{algorithm}




The algorithm for extracting frames from the video source is shown in Algorithm~\ref{alg:frame_extraction}. This algorithm efficiently captures frames in real-time for further processing. The frames extracted here serve as structured input to downstream modules, such as crowd analysis, captioning, and anomaly detection.

\section{\MakeUppercase{Crowd Estimation}}

This module estimates the number of individuals in each extracted video frame using a deep learning-based object detection model. A fine-tuned YOLO (You Only Look Once) model, trained specifically to detect the “person” class in dense and diverse scenes, is employed for accurate and real-time crowd detection.

\subsection*{Fine-Tuned YOLO Model for Person Detection}

The YOLO architecture enables fast and accurate object detection by predicting bounding boxes and class probabilities in a single forward pass. The model used here has been fine-tuned on a crowd-specific dataset to enhance performance under varied lighting, scale, and occlusion conditions.

\noindent\textbf{Pseudocode for Crowd Counting} 
\vspace{.2em} \\
\noindent\textbf{Input:} Video frame, fine-tuned YOLO model\\
\noindent\textbf{Output:} Estimated number of people in the frame
\begin{algorithmic}[1]
% \Require Video frame, fine-tuned YOLO model
% \Ensure Estimated number of people in the frame

\Function{EstimateCrowd}{frame, yolo\_model}
    \State $detections \gets \text{yolo\_model.predict}(frame)$
    \State $person\_count \gets 0$
    \ForAll{detection in $detections$}
        \If{detection.class\_label == "person"}
            \State $person\_count \gets person\_count + 1$
        \EndIf
    \EndFor
    \State \Return $person\_count$
\EndFunction
\end{algorithmic}

\section{\MakeUppercase{Graphical Representation}}

This module offers a real-time graphical visualization of crowd data to help users monitor variations and trends interactively. It serves as an intuitive interface to interpret density metrics over time.

% \subsection*{Real-Time Visualization with Chart.js and WebSockets}

Crowd count data is continuously streamed to the web interface using WebSockets, enabling low-latency updates. Chart.js renders this data dynamically in line graphs, presenting crowd size variation across frames in real time.

\section{\MakeUppercase{Fluctuation Detection}}

Detecting sudden changes in crowd density is essential for identifying abnormal events. This module implements a threshold-based algorithm to highlight significant deviations from recent trends.

The fluctuation detection approach outlined in Algorithm \ref{alg:fluctuation_detection} works by comparing the current crowd count with the historical trend within a defined window size $N$. By calculating the mean of the last $N$ crowd counts, the algorithm determines the relative percentage change between the current count and the historical average. If this change exceeds a specified threshold $T$, it flags the event as a significant fluctuation. This method provides a robust way to identify rapid shifts in crowd dynamics, ensuring that abnormal changes are detected promptly and accurately.

\begin{algorithm}[H]
\caption{Crowd Fluctuation Detection}
\label{alg:fluctuation_detection}
\textbf{Input:} Current crowd count, list of previous counts, window size $N$, and threshold $T$\\
\textbf{Output:} Boolean indicating significant fluctuation
\begin{algorithmic}[1]
\Function{DetectFluctuation}{current\_count, previous\_counts, N, T}
    \If{$|\text{previous\_counts}| < N$}
        \State \Return \textbf{False}
    \EndIf
    \State $recent \gets$ last $N$ values from $previous\_counts$
    \State $mean \gets$ average of $recent$
    \If{$mean = 0$}
        \State \Return \textbf{False}
    \EndIf
    \State $change \gets \frac{|current\_count - mean|}{mean} \times 100$
    \State \Return $change > T$
\EndFunction
\end{algorithmic}
\end{algorithm}


This logic ensures that the system remains responsive to dynamic changes, helping in early detection of potential hazards or unusual behavior patterns.

\section{\MakeUppercase{Video Captioning}}

To enhance situational awareness and support high-level semantic analysis of surveillance footage, the system includes a video captioning module that generates descriptive textual summaries of visual content. This process transforms visual data into natural language, providing a human-readable interpretation of activities and objects within a scene.

\subsection{Caption Generation Using a Vision-Language Model}

The module leverages a pre-trained Vision-Language Model (VLM), which is capable of understanding both visual inputs and their corresponding textual representations. These models are trained on large-scale datasets containing images or video frames alongside textual annotations, enabling them to associate visual patterns with natural language descriptions.

During inference, frames extracted from the video stream are passed to the VLM. The model analyzes each frame and outputs captions that describe the observed scene, such as the actions being performed, the objects present, and the relationships among them. This semantic understanding is especially useful in surveillance contexts where human operators may need quick summaries without watching full video feeds.

\subsection{Frame Sampling Strategy}

To balance computational efficiency and contextual completeness, the system implements a frame sampling mechanism. Rather than captioning every frame, it selects representative frames from each time segment (e.g., every few seconds). These sampled frames are used as input to the captioning model, ensuring coverage of important events while minimizing redundancy.

\subsection{Applications of Captions}

The generated captions serve as a foundation for downstream tasks such as anomaly detection, event summarization, and natural language-based querying. By converting raw video data into structured textual data, the system becomes more interpretable, searchable, and actionable.

\section{\MakeUppercase{Anomaly Detection}}

The anomaly detection module is designed to identify unusual or suspicious activities in surveillance footage based on semantic cues extracted from video captions. It combines natural language processing techniques with supervised learning to detect deviations from normal behavioral patterns.

\subsection{Two-Stage Pipeline for Anomaly Detection}

The anomaly detection process is implemented in two primary stages:

\begin{enumerate}
\item \textbf{Caption Summarization:}  
Since the video captioning module generates multiple descriptions over time, these captions are passed through a summarization model to condense them into a short, coherent summary that captures the essence of the scene. A sequence-to-sequence transformer model is used, trained to generate concise summaries while maintaining semantic integrity.

\item \textbf{Text Classification:}  
The summarized text is analyzed by a fine-tuned classification model for anomaly detection. It determines whether the description indicates normal or abnormal behavior, returning a label or confidence score. Transformer models like BERT have shown strong performance in natural language classification tasks due to their attention mechanisms and contextual encoding \cite{meem2023improving}, making them well-suited for anomaly detection.


\end{enumerate}

\subsection{Training and Adaptation}

Although pre-trained models are used, they are further fine-tuned on domain-specific data to adapt them to the context of surveillance footage. The classification model is trained on labeled caption summaries derived from both normal scenarios (e.g., walking, standing) and anomalous scenarios (e.g., fighting, running in restricted areas). This training enables the model to learn subtle linguistic patterns that may signify an abnormal event.

\subsection{Advantages of Text-Based Anomaly Detection}

By analyzing natural language summaries rather than raw pixel data, the system introduces an additional layer of abstraction and explainability. It allows for more transparent decision-making, as each detected anomaly is associated with a human-readable textual rationale. This approach is especially valuable in scenarios where operators need not only alerts but also contextual explanations of what went wrong.

Furthermore, using text-based models reduces computational overhead compared to dense frame-by-frame video analysis, making the solution scalable and efficient for deployment in real-time surveillance environments.

\section{\MakeUppercase{Sparse Zone Suggestion}}

This module identifies underutilized areas by analyzing spatial distributions of detected individuals.

The sparse zone suggestion process, as described in Algorithm \ref{alg:sparse_zone_detection}, helps in identifying areas with low crowd density by analyzing spatial distributions. Each zone is evaluated based on the number of detections it contains, considering the area of the zone. The algorithm computes the density of individuals in each zone and compares it to a predefined threshold $S$. If the density falls below the threshold, the zone is classified as sparse. This technique enables the system to recognize underutilized spaces, potentially assisting in optimizing crowd distribution or directing individuals to less crowded areas.


\begin{algorithm}[H]
\caption{Sparse Zone Detection}
\label{alg:sparse_zone_detection}
\textbf{Input:} List of spatial zones, current frame detections, threshold $S$\\
\textbf{Output:} List of sparse zones
\begin{algorithmic}[1]
\Function{DetectSparseZones}{zones, detections, S}
    \State $sparse\_zones \gets [\hspace{0.2em}]$
    \ForAll{zone in $zones$}
        \State $count \gets 0$
        \State $area \gets \text{CalculateArea}(zone)$
        \ForAll{detection in $detections$}
            \If{$\text{IsWithinZone}(detection.bounding\_box, zone)$}
                \State $count \gets count + 1$
            \EndIf
        \EndFor
        \State $density \gets count / area$
        \If{$density < S$}
            \State Append $zone$ to $sparse\_zones$
        \EndIf
    \EndFor
    \State \Return $sparse\_zones$
\EndFunction

\Function{IsWithinZone}{bounding\_box, zone}
    \State \Return \textbf{True} if bounding box intersects zone
\EndFunction
\end{algorithmic}
\end{algorithm}

\section{\MakeUppercase{Voice Alert System}}

This module delivers auditory notifications when specific events, such as anomalies or fluctuations, are detected. Textual alerts are sent to a Text-to-Speech (TTS) API, which returns an audio stream. The system plays this stream through speakers for hands-free real-time notification.

\section{\MakeUppercase{Optimized Routing}}

This module provides intelligent route suggestions based on both real-time and predicted crowd density information, enhancing navigation in dynamic environments.

\subsection{LSTM for Crowd Density Prediction}
Long Short-Term Memory (LSTM) networks are employed to predict future crowd densities using historical movement and congestion data. The model is structured with an input lag of 12 time steps and utilizes a deep architecture with two LSTM layers of 64 units each, followed by a dense output layer. The network is trained using the RMSprop optimizer with a Mean Squared Error (MSE) loss function. Key training hyperparameters include a batch size of 128 and 10 training epochs, with early stopping and validation split mechanisms for performance stabilization. The system evaluates the model using standard metrics such as RMSE, MAE, MAPE, and $R^2$, along with a custom threshold-based accuracy metric to assess predictions within ±10\% of actual values. Visualizations and CSV logs of loss, error metrics, and prediction accuracy aid in comprehensive performance monitoring. The predicted crowd densities help optimize navigation by guiding users toward less congested and more efficient routes.

The use of LSTM for dynamic crowd prediction is inspired by its proven effectiveness in traffic management systems. For example, Li et al. \cite{li2020deep} proposed an LSTM-based adaptive traffic controller that achieved significant reductions in delays and congestion by using spatial-temporal state representations to guide decisions under dynamic conditions.


\subsection{Path Calculation using A* Algorithm}

To determine the optimal route, the system employs the A* search algorithm, which balances the actual travel cost and an estimated cost to the destination. By integrating LSTM-based congestion predictions into the cost function, the algorithm dynamically avoids predicted high-traffic zones, improving both travel time and user experience.


The A* algorithm, as shown in Algorithm \ref{alg:a_star}, finds the optimal route by balancing travel cost and heuristic estimates. By integrating crowd density into the edge cost through LSTM-based predictions, it avoids congested areas. This adaptive approach improves both travel time and user experience. It ensures the system reacts to real-time traffic conditions efficiently.

\begin{algorithm}[H]
\caption{A* Algorithm with Density Cost}
\label{alg:a_star}
\textbf{Input:} Start node, goal node, graph with crowd density as edge cost\\
\textbf{Output:} Optimal path from start to goal
\begin{algorithmic}[1]
\Function{AStar}{start, goal, graph}
    \State $open \gets$ priority queue with $(start, 0)$
    \State $came\_from \gets$ empty map
    \State $g\_score[start] \gets 0$
    \While{$open$ is not empty}
        \State $current \gets$ node in $open$ with lowest total cost
        \If{$current == goal$}
            \State \Return \text{ReconstructPath}$(came\_from, current)$
        \EndIf
        \ForAll{neighbor in $graph.neighbors(current)$}
            \State $temp\_g \gets g\_score[current] + \text{EdgeCost}(current, neighbor)$
            \If{$temp\_g < g\_score[neighbor]$}
                \State $came\_from[neighbor] \gets current$
                \State $g\_score[neighbor] \gets temp\_g$
                \State $f\_score[neighbor] \gets temp\_g + \text{Heuristic}(neighbor, goal)$
                \State Add $neighbor$ to $open$
            \EndIf
        \EndFor
    \EndWhile
    \State \Return \textbf{Failure}
\EndFunction
\end{algorithmic}
\end{algorithm}

\subsection{User Interface via Tkinter}

The system includes a user-friendly graphical interface developed using Tkinter. Users can input their starting point and destination through this interface. Once submitted, the system processes this information and computes the best available route based on current and forecasted conditions.

\subsection{Interactive Map Visualization with Folium}

After route computation, the path is visualized on an interactive map using the Folium library. This map allows users to explore their route in a real-world context, with features like zooming, panning, and location markers, making the system more intuitive and visually informative.



\section{\MakeUppercase{Web Interface}}

\textbf{Frontend and Backend System:} The system features a reactive frontend built with SvelteJS, providing an interactive user experience with live video streams, analytical graphs, anomaly alerts, and sparse zone highlights. WebSocket communication ensures low-latency updates between the frontend and backend. Node.js serves the Svelte application, manages WebSocket connections, and exposes RESTful APIs to interact with the Python-based AI modules.

\textbf{Python–Node Integration:} Core AI models, including detection, captioning, and anomaly classification, are executed in Python. Node.js and Python are integrated through HTTP APIs, allowing the backend to trigger inference tasks and retrieve results as needed.



\vspace{1em}

This chapter outlined the implementation of all functional modules of the intelligent surveillance system. The following chapter discusses the experimental setup and evaluates the performance of the proposed system.
